\anglais
\chapter*{Abstract}

Raman spectroscopy is routinely used to characterize disorder in graphene, yet, to our knowledge, no resonant Raman spectrum of a material with defects has been computed entirely from first principles: existing models parametrize the electron-defect coupling. This thesis computes, for a vacancy in graphene, the \textit{ab initio} ingredients of the double-resonant Raman amplitudes, except for the electron-photon coupling. The intensities themselves are left for future work.

After formulating the theory of resonant Raman scattering with Feynman diagrams, we develop the central contribution of this thesis: the calculation of the electron-defect coupling, in a plane-wave basis with supercells treated with density functional theory. Zero-padding the defect potential decouples the supercell size from the Wannier-function grid and allows the coupling to be interpolated onto arbitrarily dense grids. The $T$-matrix formalism treats the scattering induced by the defect at all orders in its potential, and the localization of the coupling in the Wannier basis allows this matrix to be solved in a low-dimensional local subspace: the inversion shrinks from dimension $288\,000$ to $145$ on a $240\times 240$ grid. For an unrelaxed vacancy, the coupling is dominated by its local part, the defect-induced scattering rate reveals a quasi-bound state $0.18$~eV below the Dirac point that makes the vacancy a resonant scatterer, and the Born approximation overestimates this rate by a median factor of $46$. Finally, the phonon-induced scattering rate, computed with \texttt{EPW} in the same Wannier subspace, is of the same order of magnitude as the defect-induced one at $300$~K and a vacancy concentration of $1$\%.

{\bf Keywords:} Graphene, resonant Raman spectroscopy, density functional theory, electron-phonon coupling, electron-defect coupling, Wannier functions, density functional perturbation theory, \textit{ab initio} calculations, condensed matter physics.